\documentclass[runningheads]{llncs}

\usepackage{graphicx}
\usepackage{booktabs}
\usepackage{amsmath}
\usepackage{hyperref}
\usepackage{url}
\usepackage{microtype}
\usepackage{multirow}

\begin{document}

\title{Learning to Estimate PostgreSQL Query Cost, Planning Time, and Execution Time}

\author{%
  Don Tay Xi Cheng, Justin Leng, Lee Poe Aik, Yap Ho Wen
}
\institute{%
  National University of Singapore
}

\maketitle

%------------------------------------------------------------------------------
\begin{abstract}
PostgreSQL's built-in cost estimator is poorly calibrated to actual
execution time.
We train classical ML models on 3,300 TPC-H observations to predict three
targets: planner cost, planning time, and execution time.
On a random split, our best model achieves $R^2 = 0.958$ and median
q-error $1.056\times$ for execution time, versus $382.6\times$ for
PostgreSQL's own cost proxy.
Under 5-fold template-grouped cross-validation, q-error rises to
$2.29\times$, which exposes a generalisation gap between seen and unseen
query structures.
Plan features dominate; SQL features add marginal value.
Planning time is the hardest target ($R^2 = 0.45$) due to measurement
noise at sub-millisecond granularity.
\end{abstract}

%------------------------------------------------------------------------------
\section{Introduction}

PostgreSQL selects query execution plans by minimising a heuristic cost
estimate.
This cost is unitless and intended for plan ranking, not runtime
prediction~\cite{postgresql}.
Using it directly as a latency predictor yields median q-error
$382.6\times$ in our experiments.

An alternative is to train regressors on queries executed with
\texttt{EXPLAIN (ANALYZE, FORMAT JSON)} and predict performance from
query and plan features directly.
Prior work achieves strong results~\cite{marcus2019,zhao2022} but
typically uses deep learning on tens of thousands of queries and focuses
on a single target.

This paper contributes: (i) a reproducible TPC-H collection and feature
extraction pipeline producing 3,300 labelled observations; (ii) a
comparative evaluation of four model families on three targets jointly;
(iii) feature ablations separating SQL from plan signal; and (iv) a
5-fold template-grouped cross-validation protocol that quantifies the
generalisation gap on structurally unseen queries.

%------------------------------------------------------------------------------
\section{Background}
\label{sec:background}

\subsection{Query Optimisation in PostgreSQL}

PostgreSQL's cost-based optimiser selects the plan with the lowest
estimated cost, derived from table statistics and operator cost constants.
Running \texttt{EXPLAIN (ANALYZE, BUFFERS, FORMAT JSON)} exposes three
targets: (i)~\emph{planner total cost} (unitless, used for plan
selection), (ii)~\emph{planning time} (ms, time to generate the plan),
and (iii)~\emph{execution time} (ms, actual wall-clock query runtime).
These three targets differ in scale, variance, and how well they can be
predicted from structural features.

\subsection{TPC-H Benchmark}

TPC-H~\cite{tpch} is a standard decision-support benchmark with 22
parameterised query templates over an 8-table schema covering joins,
aggregations, and range predicates.
Query instances are generated by \texttt{QGEN}; the scale factor (SF)
controls data volume.
We use SF~1.0 (${\approx}$1\,GB), consistent with prior
work~\cite{marcus2019,liang2023}.

\subsection{Evaluation Metrics}

\paragraph{Q-error~\cite{leis2015}.}
$\text{q-error} = \max(\hat{y}/y,\; y/\hat{y}) \geq 1$.
A value of 1.0 is perfect; 2.0 means the estimate is off by $2\times$ in
either direction.
Q-error is standard in cost estimation because multiplicative errors
matter more than additive ones.
We report it at the 50th, 90th, 95th, and 99th percentiles.

\paragraph{RMSE, MAE, MAPE, $R^2$.}
Root mean squared error and mean absolute error report absolute error in
the target's native unit.
MAPE and $R^2$ report relative error and explained variance respectively.

%------------------------------------------------------------------------------
\section{Related Work}
\label{sec:related}

Leis et al.~\cite{leis2015} established q-error as the metric for
cardinality estimation and showed PostgreSQL's estimator fails on
multi-join queries with correlated predicates.
BICE~\cite{liang2023} trains a bidirectional ensemble on 22,000 TPC-H
queries for cost and cardinality estimation; our dataset is smaller but
fully reproducible.
QPPNet~\cite{marcus2019} encodes plan trees with a recursive neural
network to predict latency on TPC-H and TPC-DS at tens of thousands of
queries.
QueryFormer~\cite{zhao2022} applies a plan-tree transformer and notes
TPC-H is ``comparatively simple'' relative to real workloads, motivating
our template-grouped generalisation test.

Our work differs in three ways: we predict three targets jointly rather
than execution time alone, we use classical ML suited to our dataset
size, and we measure the template-grouped generalisation gap, which
prior TPC-H studies typically omit.

%------------------------------------------------------------------------------
\section{Methodology}
\label{sec:method}

\subsection{Data Collection}

We provision PostgreSQL~16 in Docker and load TPC-H at SF~1.0.
For each of the 22 templates we generate 50 parameter sets
deterministically (seed 20260322), giving 1,100 query instances each
executed three times with \texttt{EXPLAIN (ANALYZE, BUFFERS, FORMAT JSON)}.
The median of three warm-cache runs is used as the execution-time label.
Timeout is 300\,s; all 3,300 observations at SF~1.0 completed
successfully.
Each record stores the raw SQL, JSON plan, planner cost, planning time,
execution time, template ID, parameter set ID, and scale factor.

\subsection{Feature Extraction}

\paragraph{SQL features (8).}
Parsed from each SQL string using SQLGlot~\cite{sqlglot}: join count,
predicate count, aggregation presence, group-by count, order-by count,
limit count, subquery count, and selected column count.

\paragraph{Plan features (15+).}
Extracted by recursively traversing the PostgreSQL JSON plan: total node
count, join/scan/aggregate/sort node counts, max tree depth, estimated
row counts (sum and max), output widths (sum and max), and startup/total
cost aggregates.

\paragraph{Null handling.}
A binary null-indicator is appended per feature; null values are filled
with zero.

\subsection{Models and Training}

We compare four families: \textbf{Dummy Mean} (unconditional mean
baseline), \textbf{Ridge} ($L_2$ regression with \texttt{StandardScaler}),
\textbf{Random Forest} (200 trees, \texttt{min\_samples\_leaf=5}), and
\textbf{Histogram Gradient Boosting}~\cite{sklearn}.
Execution time is log-transformed (\texttt{log1p}) before fitting to
handle right-skewed runtimes.
The family with the best validation RMSE is selected per target and
refit on train\,+\,validation before test evaluation.

\subsection{Evaluation Protocol}

\paragraph{Random split.}
A 60/20/20 train/val/test partition at the query-instance level measures
in-distribution performance.

\paragraph{5-fold template-grouped CV.}
The 22 templates are divided into 5 folds (${\approx}$4 templates each).
Each fold trains on the remaining 18 templates and tests on the 4
held-out templates; train and test sets are strictly disjoint.
Metrics are averaged across folds.
This is strictly harder than the random split and measures generalisation
to unseen query structures.

\paragraph{Ablations.}
We compare SQL-only, plan-only, and combined features using Histogram
Gradient Boosting on the random split, and additionally evaluate a
\emph{PostgreSQL cost proxy} (planner \texttt{total\_cost} used directly
as the execution-time prediction).

%------------------------------------------------------------------------------
\section{Performance Evaluation}
\label{sec:eval}

\noindent\textbf{Setup.}
MacBook (Apple Silicon), macOS~25, Python~3.14, scikit-learn, polars,
SQLGlot.
The full pipeline reproduces from \texttt{uv run ivory train baseline \&\&
uv run ivory evaluate}.

\subsection{Baseline Model Comparison}

Table~\ref{tab:baseline} and Figure~\ref{fig:model_comparison} summarise
test-set results.
Random Forest is selected for execution time ($R^2=0.958$,
q-err\,p50\,$=1.056\times$); Histogram Gradient Boosting for planner cost
($R^2=0.999$) and planning time ($R^2=0.448$).
Ridge and Dummy Mean lag far behind; non-linear models are necessary for
these targets.

\begin{table}[t]
\centering
\caption{Test-set metrics for all model families. \textbf{Bold} = selected
model per target. RF = Random Forest, HGB = Hist.\ Gradient Boosting.}
\label{tab:baseline}
\setlength{\tabcolsep}{5pt}
\begin{tabular}{llrrrr}
\toprule
\textbf{Target} & \textbf{Model} & \textbf{RMSE} & \textbf{MAE} &
\textbf{Q-err\,p50} & \textbf{$R^2$} \\
\midrule
\multirow{4}{*}{\shortstack[l]{Exec.\ time\\(ms)}}
  & Dummy Mean  & 484.85 & 288.04 & 1.748 & $-$0.168 \\
  & Ridge       & 237.20 & 148.53 & 1.321 & 0.721 \\
  & \textbf{RF} & \textbf{91.99} & \textbf{38.13} & \textbf{1.056} & \textbf{0.958} \\
  & HGB         & 90.00  & 37.21  & 1.056 & 0.960 \\
\midrule
\multirow{4}{*}{\shortstack[l]{Planner\\cost}}
  & Dummy Mean   & 140,727 & 84,649 & 1.496 & $-$0.009 \\
  & Ridge        & 78,376  & 63,171 & 1.569 & 0.687 \\
  & RF           & 4,239   & 1,144  & 1.001 & 0.999 \\
  & \textbf{HGB} & \textbf{4,301} & \textbf{1,290} & \textbf{1.001} & \textbf{0.999} \\
\midrule
\multirow{4}{*}{\shortstack[l]{Planning\\time (ms)}}
  & Dummy Mean   & 1.246 & 0.813 & 1.703 & $-$0.001 \\
  & Ridge        & 0.942 & 0.549 & 1.440 & 0.427 \\
  & RF           & 0.922 & 0.498 & 1.289 & 0.452 \\
  & \textbf{HGB} & \textbf{0.925} & \textbf{0.499} & \textbf{1.279} & \textbf{0.448} \\
\bottomrule
\end{tabular}
\end{table}

\begin{figure}[t]
  \centering
  \includegraphics[width=\textwidth]{fig_model_comparison.png}
  \caption{Model family comparison by target. Left: RMSE. Right: Q-error p50.
  Bold bars indicate the selected model per target.}
  \label{fig:model_comparison}
\end{figure}

Execution time is predicted to within 6\% at the median
(q-err\,p50\,$=1.056\times$, RMSE\,=\,91.99\,ms, $R^2=0.958$).
Planner cost is near-perfectly predicted ($R^2=0.999$,
q-err\,p50\,$=1.001\times$), as expected: it is a deterministic function
of plan structure that plan features directly encode.
Planning time is the hardest target ($R^2=0.448$); PostgreSQL planning
takes under 5\,ms, where CPU jitter dominates over structural signal.

\subsection{Generalisation Under Template-Grouped CV}

\begin{table}[t]
\centering
\caption{Random-split vs.\ 5-fold template-grouped cross-validation results.}
\label{tab:grouped}
\begin{tabular}{lrrrr}
\toprule
\textbf{Target} &
\multicolumn{2}{c}{\textbf{Random split}} &
\multicolumn{2}{c}{\textbf{Grouped (5-fold avg.)}} \\
\cmidrule(lr){2-3}\cmidrule(lr){4-5}
 & Q-err\,p50 & Q-err\,p90 & Q-err\,p50 & Q-err\,p90 \\
\midrule
Exec.\ time (ms) & 1.056 & 1.280 & 2.291 & 3.324 \\
Planner cost     & 1.001 & 1.060 & 1.880 & 6.135 \\
Planning time (ms) & 1.279 & 2.416 & 1.937 & 3.374 \\
\bottomrule
\end{tabular}
\end{table}

\begin{figure}[t]
  \centering
  \includegraphics[width=\textwidth]{fig_random_vs_grouped.png}
  \caption{Q-error p50 and p90 under random split vs.\ template-grouped CV.}
  \label{fig:grouped}
\end{figure}

Table~\ref{tab:grouped} and Figure~\ref{fig:grouped} show a generalisation
gap across all targets.
For execution time, median q-error rises from $1.056\times$ to
$2.29\times$ on unseen templates.
Planner cost degrades more sharply at the tail (q-err\,p90:
$1.06\times \to 6.13\times$); certain plan structures are particularly
difficult to extrapolate to.
Random-split performance is therefore optimistic: the model partially
memorises template-level patterns.

\subsection{Feature Ablation}

\begin{table}[t]
\centering
\caption{Feature ablation results (Histogram Gradient Boosting, random
split). ``PG proxy'' uses \texttt{total\_cost} directly as execution-time
prediction.}
\label{tab:ablation}
\setlength{\tabcolsep}{5pt}
\begin{tabular}{llrrrr}
\toprule
\textbf{Target} & \textbf{Condition} & \textbf{RMSE} & \textbf{MAE} &
\textbf{Q-err\,p50} & \textbf{$R^2$} \\
\midrule
\multirow{4}{*}{\shortstack[l]{Exec.\ time\\(ms)}}
  & SQL only  & 97.75   & 37.03   & 1.054 & 0.953 \\
  & Plan only & 88.89   & 36.86   & 1.055 & 0.961 \\
  & Combined  & 89.99   & 37.21   & 1.056 & 0.960 \\
  & PG proxy  & 207,202 & 152,988 & 382.6 & $-$213,290 \\
\midrule
\multirow{3}{*}{\shortstack[l]{Planner\\cost}}
  & SQL only  & 6,863 & 1,843 & 1.001 & 0.998 \\
  & Plan only & 4,304 & 1,372 & 1.002 & 0.999 \\
  & Combined  & 4,301 & 1,290 & 1.001 & 0.999 \\
\midrule
\multirow{3}{*}{\shortstack[l]{Planning\\time (ms)}}
  & SQL only  & 0.918 & 0.492 & 1.263 & 0.457 \\
  & Plan only & 0.924 & 0.498 & 1.277 & 0.449 \\
  & Combined  & 0.925 & 0.499 & 1.279 & 0.448 \\
\bottomrule
\end{tabular}
\end{table}

\begin{figure}[t]
  \centering
  \includegraphics[width=\textwidth]{fig_ablation_execution_time.png}
  \caption{Feature ablation for execution time. The dashed red line marks
  the PostgreSQL cost proxy q-err\,p50 ($382.6\times$).}
  \label{fig:ablation}
\end{figure}

Plan features dominate: plan-only achieves RMSE\,=\,88.89 vs.\ 97.75 for
SQL-only (9.1\% better).
Adding SQL features to the plan produces no gain---combined RMSE\,=\,89.99
is marginally worse than plan-only, consistent with SQL features being
largely redundant once the plan is available: the plan already encodes
join structure, predicate selectivity estimates, and aggregation layout.
For planning time, all three feature subsets perform identically
($R^2 \approx 0.45$): the ceiling is due to target noise, not feature
inadequacy.
The PostgreSQL cost proxy (q-err\,p50\,$=382.6\times$) is catastrophically
worse than even the weakest learned model (SQL-only,
q-err\,p50\,$=1.054\times$)---a three-order-of-magnitude gap between planner
cost and actual runtime.

\subsection{Error Analysis by Template}

\begin{figure}[t]
  \centering
  \includegraphics[width=\textwidth]{fig_q_error_by_template.png}
  \caption{Q-error distribution per template (execution time, fold 1),
  sorted by median q-error descending.}
  \label{fig:template_error}
\end{figure}

Figure~\ref{fig:template_error} shows Q18 is the hardest template (median
q-error $2.94\times$, max $5.45\times$).
Q18 contains a \texttt{HAVING COUNT(*) > 300} predicate over a large
hash aggregation; its execution time varies sharply with parameter values
in ways 50 training samples cannot capture.
Q3 is moderately hard ($1.83\times$); Q4 and Q13 are easier
($1.24\times$ and $1.68\times$ respectively), though Q13 has a longer
tail (max $3.66\times$).
Hard templates fail consistently across folds, not randomly.

\subsection{Discussion}

Plan features subsume SQL features because the plan is produced after SQL
parsing and already encodes join depth, aggregation structure, and
predicate complexity.
For planning time, the low $R^2$ reflects measurement noise: at
sub-millisecond scales, CPU scheduling and cache state add more variance
than query structure does, and no feature set can close this gap.
The Q18 difficulty is different in kind. Its \texttt{HAVING COUNT(*) > 300}
clause makes execution time depend on how many groups survive the
predicate, which varies non-linearly with parameter values and cannot be
recovered from plan structure alone.

%------------------------------------------------------------------------------
\section{Conclusion}
\label{sec:conclusion}

We trained and evaluated ML models to predict three PostgreSQL query
performance targets on TPC-H.
Key findings: (i) learned models reduce execution-time error by three
orders of magnitude over PostgreSQL's cost proxy (q-err\,p50:
$1.056\times$ vs.\ $382.6\times$); (ii) planner cost is near-perfectly
predicted ($R^2=0.999$), planning time is noise-limited ($R^2=0.45$);
(iii) plan features alone provide essentially all signal; and (iv)
template-grouped CV ($2.29\times$) better reflects expected accuracy on
unseen query structures than the random split ($1.056\times$).

Limitations include a single scale factor (SF~1.0) and TPC-H's 22
fixed templates, which are structurally simpler than real
workloads~\cite{leis2015,zhao2022}.
Future work includes multi-scale training, the TPC-DS workload (99
templates), tree-structured encoders~\cite{marcus2019}, and using
predicted cost to guide query reformulation.

%------------------------------------------------------------------------------
\bibliographystyle{splncs04}
\bibliography{references}

\end{document}
